\section{課題解決のための提案}

前章では、eスポーツの市場規模やファン人口が増加している一方で、
地域やゲームコミュニティなどが開催する大会では、競技者情報や試合結果の管理、
情報共有に手作業が発生しやすいことなどを課題として挙げました。また、eスポーツ大会では、
競技者だけではなく観戦者に対しても、大会の進行状況や試合結果などの情報を提供することが重要です。
特に、注目選手が参加する大会では、その選手の試合状況や今後の対戦予定を確認したい観戦者も存在します。

そこで、本システムでは大会運営に必要な情報を一元管理し、主催者・競技者・観戦者がそれぞれ必要な情報を確認できる
「小規模ゲーム大会運営システム(Play Line)」を提案します。本システムでは、競技者の参加登録から
試合結果の管理までを効率化するとともに、大会情報とライブ配信へのアクセスを関連付けることで、
競技者が自分の試合情報を確認しやすくし、観戦者が大会全体や注目選手の試合を追いかけやすい環境の実現を目指します。

\subsection{小規模ゲーム大会運営システム(Play Line)}

本システムでは、大会の参加登録から大会終了までに必要となる競技者情報、対戦情報、試合結果などを
データベースで一元管理します。これにより、DiscordやExcel、Googleスプレッドシートなどに分散していた
大会日程、参加者情報、対戦表、試合結果などを一元管理し、主催者の運営負担を軽減するとともに、
参加者や観戦者が必要な情報を容易に確認できるようにします。

まず、主催者が本システムに大会情報や参加条件などを登録し、競技者の募集を行います。
競技者はWebブラウザから本システムのWebサイトへアクセスし、参加したい大会への登録を行います。
登録された情報はデータベースに保存され、主催者は登録された競技者の情報を一覧で確認できるようにします。

大会開催時には、登録された競技者情報を利用してトーナメント表を作成します。
競技者はWebサイトから自分の対戦相手やトーナメント上の位置を確認できます。
また、主催者は試合終了後に試合結果を登録し、その結果に基づいて大会の進行状況を更新します。
これにより、競技者は次の対戦相手や次の試合を確認できるようになります。
主催者が登録・更新した大会情報や試合結果はデータベースに保存され、観戦者にも提供されます。
観戦者は大会ページから、トーナメント表、試合結果、現在行われている試合、今後行われる試合などを確認できます。

本システムにおける主催者による情報の登録・更新から、競技者や観戦者への情報提供までの流れを以下に示します。

\paragraph{主催者}
\begin{itemize}
  \item 大会情報、参加条件の登録
  \item 競技者情報の確認
  \item トーナメント表の作成
  \item 試合結果の登録、変更、大会情報の更新
\end{itemize}
\paragraph{競技者}
\begin{itemize}
  \item 大会への参加登録
  \item トーナメント表における現在の位置の確認
  \item 試合結果の登録、次の対戦相手の確認
\end{itemize}
\paragraph{観戦者}
\begin{itemize}
  \item 大会全体のトーナメント表
  \item 試合結果の確認
  \item 現在行われている試合、今後予定されている試合の確認
  \item 注目選手の過去の試合結果、現在の試合状況の確認
  \item 注目選手の次の対戦相手、トーナメント上の位置の確認
  \item ライブ配信URLから配信ページへのアクセス
\end{itemize}

このように、大会運営に必要な情報を一つのシステムで管理することで、主催者による
情報の整理や更新の負担を軽減します。また、主催者が登録・更新した大会情報を
競技者と観戦者が同じWebサイトから確認できるようにすることで、個別に情報を伝達する手間を減らし、
対戦相手や試合結果などの情報を効率的に共有できるようにします。
これにより、利用者が必要な情報を容易に取得できる仕組みの実現を目指します。

\subsection{注目選手の参加を活かした観戦者向け情報提供}

1章で述べたように、注目選手の参加は大会への関心を高める要因の一つとなる可能性があります。
そこで本システムでは、大会運営の効率化に加え、観戦者が関心のある選手の試合を継続して
追いやすくするための情報提供を行います。例えば、大会ページから

\begin{itemize}
  \item 注目選手（運営による注目）の現在の対戦状況
  \item 注目選手の過去の全ての試合結果
  \item 次の対戦相手
  \item トーナメント上の現在位置
  \item 大会全体の進行状況
  \item 各試合の結果
\end{itemize}
などを確認できるようにします。

これにより、特定の選手を応援している観戦者が、その選手の現在の状況や次の試合を把握しやすくなります。
また、地域大会において実績のある選手や地域で人気のある選手が参加した場合にも、
その選手の対戦状況や試合結果をシステム上で確認できるようにすることで、
会場やオンラインで大会を観戦するきっかけを作り、大会への関心を高めることを目指します。

eスポーツの観戦動機に関する研究では、1,052人を対象とした調査により、好きな選手への
憧れや応援等の動機が、ライブ配信の視聴行動に関係することが報告されています。\cite{scientificreports2024}
\footnote{Yuxin Tang, Xiaowen Zhang, Shengfeng Zan, “Exploring e-sports fans’ motivation for watching live streams based on self-determination theory,” Scientific Reports, Vol. 14, Article 13858, 2024. 

自己決定理論に基づくeスポーツファンのライブ配信視聴動機を探る |科学報告書　（閲覧日：2026年10月15日） }

このことから、観戦者が関心のある選手の試合情報を確認できることは、試合への関心を高め、
観戦を楽しむための一助になると考えられます。
そこで、本システムでは観戦者が選手の情報や現在の対戦状況、過去の試合結果、次の対戦相手、
トーナメント表における現在の位置などを大会ページから確認できるようにします。
これにより、観戦者が注目選手を探す手間を減らし、注目選手やお気に入りの選手の試合状況を把握しやすくすることを目指します。

\subsection{大会運営情報の一元管理}

本システムでは、大会の作成から試合結果の登録までに必要な情報を一元管理します。主催者は大会名や開催日時、参加条件などを登録し、競技者の募集を行います。競技者は登録画面から氏名やプレイヤー名、メールアドレスなどを入力して参加を申し込みます。登録された情報はデータベースに保存され、主催者は参加者の情報を一覧で確認できます。

大会開始後は、主催者がトーナメント表を作成し、各試合の結果を登録します。登録された結果に応じて、大会の進行状況や次の対戦相手などの情報を更新します。これらの情報を同じシステム内で管理することで、複数のツールを使い分ける必要を減らし、大会運営に関する情報の確認や更新を容易にします。

\subsection{競技者向けの情報提供}
競技者向けには、自分の試合に関する情報を確認できるページを提供します。具体的には、自分の対戦相手や試合日程、トーナメント表における現在の位置、試合結果などを確認できるようにします。また、試合結果が登録された際には、その結果に応じて次の対戦相手や大会の進行状況も確認できます。

これにより、競技者は大会の進行状況を把握するために、主催者へ個別に問い合わせたり、複数のツールを確認したりする手間を減らすことができます。
\subsection{観戦者向けの情報提供}
観戦者向けには、大会全体の状況や注目選手に関する情報を確認できるページを提供します。大会全体のトーナメント表や試合結果、現在進行中の試合、今後予定されている試合などを表示することで、観戦者が大会の進行状況を把握できるようにします。

また、本システムでは、観戦者が特定の選手に注目して大会を追えるように、注目選手に関する情報を提供します。注目選手の過去の試合結果や現在の試合状況、次の対戦相手、トーナメント表における位置などを確認できるようにすることで、選手の大会での活躍や今後の試合を追いやすくします。

さらに、主催者が大会情報にライブ配信URLを登録できるようにし、観戦者向けのページから配信ページへアクセスできるようにします。これにより、観戦者は大会の進行状況や注目選手の情報を確認しながら、試合の配信も視聴できます。


以上の提案により、本システムでは、大会運営に必要な情報を一元的に管理することで
主催者の作業負担を軽減し、参加者には自分の試合に必要な情報を提供し、
観戦者には大会全体や注目選手の試合を追いやすい環境を提供します。

これにより、小規模なゲーム大会を単に「ゲームをプレイする場」とするのではなく、
主催者が運営しやすく、参加者が競技しやすく、観戦者も大会を楽しめる場として発展させることを目指します。
